\section{Methodology: ARS Exp D (Streaming)} \label{sec:streaming}
ARS Exp D introduces the first architecture in the phylogeny that is simultaneously parallel and streamable. It achieves this by shifting the computational cost from the final aggregation phase to a concurrent background pipeline using a \textbf{Spatial Run Registry}.

\subsection{Epoch-Based Micro-Batching}
To maintain the lock-free ingestion required for high-velocity data streams, ARS collects incoming elements in thread-local buffers called \textit{Epochs}. Once an epoch reaches its capacity threshold (typically $N_e = 32,768$), it is atomically dispatched to a background work queue \cite{zahariaDiscretizedStreamsFaulttolerant2013}. This micro-batching strategy preserves the sequential memory access patterns required for high cache locality while ensuring that the main ingestion thread is never blocked by sorting operations.

To prevent ingestion overflow, Exp D utilizes a bounded pipeline capacity. When the background work queue reaches saturation, the ingestion thread exerts back-pressure, ensuring that memory consumption remains bounded even at extreme scales ($N=10^8$).

\subsection{Concurrent Spatial Consolidation}
The core innovation of Exp D is the background consolidation process, which prevents the unbounded growth of isolated runs, conceptually similar to LSM-tree compaction \cite{oneilLogstructuredMergetreeLSMtree1996}. A background worker pool continuously monitors the work queue and executes the logic described in Algorithm \ref{alg:consolidation}.

\begin{algorithm}
\caption{Concurrent Spatial Consolidation}
\label{alg:consolidation}
\begin{algorithmic}[1]
\REQUIRE Work Queue $Q$, Spatial Registry $R$
\LOOP
    \STATE $epoch \leftarrow Q.pop()$
    \STATE $SortedEpoch \leftarrow \text{ARS\_Apex}(epoch)$
    \STATE $newRun \leftarrow \text{SpatialRun}(SortedEpoch)$
    \STATE $merged \leftarrow \text{false}$
    \FORALL{$run \in R$}
        \IF{$(run.min \gg 48) = (newRun.min \gg 48)$}
            \STATE $R.replace(run, \text{merge}(run, newRun))$
            \STATE $merged \leftarrow \text{true}$
            \STATE \textbf{break}
        \ENDIF
    \ENDFOR
    \IF{$\neg merged$}
        \STATE $R.add(newRun)$
    \ENDIF
\ENDLOOP
\end{algorithmic}
\end{algorithm}

The registry utilizes a \textbf{Bit-Mask Proximity Check} to identify spatially adjacent runs. By comparing the most significant bits (e.g., bits 48--63) of the spatial keys, the algorithm groups runs that belong to the same "spatial neighborhood." When an adjacency is detected, the background worker performs an $O(N)$ linear merge of the two runs. This proactive consolidation ensures that the number of distinct runs remains tightly bounded, reducing the final aggregation latency from an unpredictable multi-way merge to a highly bounded, cache-efficient hierarchical merge.
